\documentclass[10pt,largemargins]{homework}
\usepackage[utf8]{inputenc}
\usepackage{amsmath}
\usepackage{tikz}


\newcommand{\hwname}{ENSA-Safi}
\newcommand{\hwemail}{}
\newcommand{\hwtype}{Travaux Dirigés}
\newcommand{\hwnum}{1}
\newcommand{\hwclass}{Probabilité}
\newcommand{\hwlecture}{1}
\newcommand{\hwsection}{Z}




\begin{document}
\maketitle

%{{{ Question arbe
\question*{Modèle arbre}

Omar vérifie l'état de la météo. Il elle est bonne, il part pour une marche.
Par contre, si elle pleut, Omar va soit rester chez lui, soit sortir. En cas
il sort, il se peut qu'il se rappelle de ramener sa Parapluie. 
\begin{figure}[htpb]
    \centering
    \begin{tikzpicture}[xscale=2.5,yscale=.8]
        
        \coordinate (A) at (0,0); 
        \coordinate (B) at (1,1); 
        \node at (1.1,1) {$1$}; 
        \coordinate (C) at (1,-1); 
        \coordinate (D) at (2,0); 
        \node at (2.1,0) {$2$}; 
        \coordinate (E) at (2,-2); 
        \coordinate (F) at (3,-1); 
        \node at (3.1,-1) {$3$}; 
        \coordinate (G) at (3,-3); 
        \node at (3.1,-3) {$4$}; 

        \path[draw,thick] (A)--(B);
        \path[draw,thick] (A)--(C) -- (D);
        \path[draw,thick] (C)--(E) -- (F);
        \path[draw,thick] (C)--(E) -- (G);
    \end{tikzpicture}
    
    \caption{Modelé Arbre}
    \label{fig:omar}
\end{figure}


\begin{arabicparts}
    \small
    \item Dans le diagramme, identifier la feuille (nombre) qui correspond au cas suivant: \textbf{Elle pleut et Omar reste
        chez lui}.
    \item Ajouter une étiquette aux autre cas restants.
\end{arabicparts}
%}}}
% Question Disjoints {{{ %

\question%
Soit $A$ et $B$ deux éventent du même espace d'états $\Omega$, avec $P(A)
= 0.6$ et $P(B) = 0.7$.

\begin{arabicparts}
    \item Est ce que $A$ et $B$ peuvent être
        \textbf{disjoints}. Justifier votre réponse.
\end{arabicparts}
% }}} Question Disjoints %
% Properrite simples {{{ %
\question*{Propriétés simples}
Soient $A$, $B$ et $C$ des évènements de l'espace d'états $\Omega$. Pour chaque
propriétés, indiquer si elle est juste ou non en justifiant votre
réponse.


\begin{arabicparts}
    \item $ P(A) + P(A^c) + P(B) = P(A\cup A^c \cup C)$
    \item $ P(A) + P(B) \leq 1$
    \item $ P(A^c) + P(B) \leq 1$
    \item $ P(A\cup B \cup C) \geq P(A\cup B)$
\end{arabicparts}
% }}} Properrite simples %
% Properrite simples {{{ %
\question*{Autres propriétés}
Soient $A$, $B$ et $C$ des évènements de l'espace d'états $\Omega$. Pour chaque
propriétés, indiquer si elle est juste ou non en justifiant votre
réponse.


\begin{arabicparts}
\item $ P\left( (A\cap B) \cup ( C\cap A^c)\right) \leq P(A\cup B \cup C)$
\item $P(A\cup B\cup C) = P(A\cap C^c) + P(C) + P(B\cap A^c \cap C^c)$
\end{arabicparts}
% }}} Properrite simples %
% Tetrahedral1 {{{ %
\question*{Exemple discret}

On considère l'exemple de lance d'un dé (a quatre faces) ou tous les $16$ cas
possibles on la même probabilités. Calculer la probabilité de chaque évènement
suivant:
\begin{figure}[htpb]
\begin{center}
\begin{tikzpicture}[scale=.8, transform shape]
    \draw[thick, gray] (0,0) grid (5,5);
\end{tikzpicture}
\end{center}
\end{figure}


\begin{arabicparts}
    \item La valeur du premier dé est supérieure a la deuxième.
    \item La somme des valeurs obtenues est un nombre pair.
\end{arabicparts}


% }}} Tetrahedral1 %
%  Calcul probabilite continue {{{ %
\question*{Exemple continu}

On considère un espace d'états:

\begin{equation*}
    \Omega = [0,1]\times[0,2]:
\end{equation*}

qui contient les couples $(x,y)$ tel que $0\leq x\leq 1$ et $0\leq y\leq 2$. On
considère une loi \textbf{uniforme}. Calculer les probabilités des
évènements suivants:

\begin{arabicparts}
    \item Les deux composantes $x$ et $y$ possèdent la même valeur.
    \item La valeur de $x$, du premier lancé, est supérieure ou
        égale
        a la valeur $y$, du deuxième lance.
    \item La valeur $x^2$ est supérieure ou égale à la valeur de $y$.
\end{arabicparts}



% }}}  Calcul probabilite continue %
% Somme denombrable {{{ %
\question*{Somme dénombrable}
On considère l'espace des états $\Omega$ qui consiste des entiers naturels
$\mathbb{N}$. On suppose aussi que:

\begin{equation}
    \mathbf{P}(n) = \dfrac{1}{2^n}\quad \forall n \geq 1
\end{equation}

      Calculer la probabilité de l'évènement
        $A = \left\{3,6,9,\ldots\right\}.$. I.e, que le nombre tire
        est un \textbf{multiple} de trois.

% }}} Somme denombrable %
% Somme difference symmetrique {{{ %
\question*{Probabilité Somme symétrique}
Dans cet exercice, on se propose de prouver la formule de probabilité de
la \textbf{différence symétrique} $A\Delta B$\footnote{$A\Delta B = (A\cup B^c) \cup (B \cap
A^c)$} de deux évènements.

\begin{equation}
   P(A \Delta B) = P(A) + P(B) - 2 P(A\cap B) 
\end{equation}

% }}} Somme difference symmetrique %
% Genies {{{ %
\question*{Sport et Cinema}

Dans une classe, $\mathbf{60\%}$ pratiquent un sport, $\mathbf{70\%}$
aiment le cinéma et $\mathbf{40\%}$ sont des les deux \textbf{catégories}.

\begin{itemize}
    \item Calculer la probabilité, qu'un étudiant aléatoirement
        choisi ne pratique pas se sport et il ne n'aime pas le cinéma.
\end{itemize}

% }}} Genies %

\end{document}
